% gerado por figuras_etapa0.py em 15/11/2026 12:16: nao edite (rode o script de novo)
\newif\ifchofim \chofimfalse
\newif\ifpcparou \pcparoutrue
\newif\ifchocincopronto \chocincoprontofalse
\newif\ifchocincopiloto \chocincopilotofalse
\newif\ifchoaicincoiniciado \choaicincoiniciadofalse
\newif\ifchoaicincopronto \choaicincoprontofalse
\newif\ifchoaicincofalhas \choaicincofalhasfalse
\newif\ifchosetepronto \choseteprontofalse
\newif\ifchosetepiloto \chosetepilotofalse
\newif\ifchoaiseteiniciado \choaiseteiniciadofalse
\newif\ifchoaisetepronto \choaiseteprontofalse
\newif\ifchoaisetefalhas \choaisetefalhasfalse
\newif\ifformdezpronto \formdezprontofalse
\newif\ifformdezpassoacima \formdezpassoacimatrue
\newif\ifformvintepronto \formvinteprontofalse
\newif\ifformvintepassoacima \formvintepassoacimatrue
\newif\ifaiforminiciado \aiforminiciadofalse
\newif\ifaiformfalhas \aiformfalhasfalse
\newif\ifaiformpronto \aiformprontofalse
\newif\ifchoburaco \choburacotrue
\newcommand{\HoraDados}{15/11/2026, 12h16} % relogio do PC ao rodar o script
\newcommand{\MinWilson}{20} % ../README.md (criterios gerais)
\newcommand{\RbNDireto}{2000} % numeros_rebaixamento.tex (relatorio_rebaixamento.pdf)
\newcommand{\RbNfDireto}{314} % numeros_rebaixamento.tex (relatorio_rebaixamento.pdf)
\newcommand{\RbPfDireto}{15{,}70} % numeros_rebaixamento.tex (relatorio_rebaixamento.pdf)
\newcommand{\RbCovPfDireto}{5{,}2} % numeros_rebaixamento.tex (relatorio_rebaixamento.pdf)
\newcommand{\RbNBanda}{1000} % numeros_rebaixamento.tex (relatorio_rebaixamento.pdf)
\newcommand{\RbPfBanda}{15{,}87} % numeros_rebaixamento.tex (relatorio_rebaixamento.pdf)
\newcommand{\RbCovPfBanda}{4{,}4} % numeros_rebaixamento.tex (relatorio_rebaixamento.pdf)
\newcommand{\RbDifBanda}{0{,}16} % numeros_rebaixamento.tex (relatorio_rebaixamento.pdf)
\newcommand{\RbGanhoBanda}{2{,}7} % numeros_rebaixamento.tex (relatorio_rebaixamento.pdf)
\newcommand{\RbGanhoLinear}{4{,}0} % numeros_rebaixamento.tex (relatorio_rebaixamento.pdf)
\newcommand{\RbFracGanhoBanda}{69} % numeros_rebaixamento.tex (relatorio_rebaixamento.pdf)
\newcommand{\RbMuDoisBanda}{2{,}06} % numeros_rebaixamento.tex (relatorio_rebaixamento.pdf)
\newcommand{\RbSC}{\ensuremath{-}1{,}30} % numeros_rebaixamento.tex (relatorio_rebaixamento.pdf)
\newcommand{\RbSPhi}{\ensuremath{-}0{,}49} % numeros_rebaixamento.tex (relatorio_rebaixamento.pdf)
\newcommand{\RbSK}{\ensuremath{-}0{,}36} % numeros_rebaixamento.tex (relatorio_rebaixamento.pdf)
\newcommand{\RbKModos}{20} % numeros_rebaixamento.tex (relatorio_rebaixamento.pdf)
\newcommand{\RbNModos}{1000} % numeros_rebaixamento.tex (relatorio_rebaixamento.pdf)
\newcommand{\RbPfModos}{17{,}89} % numeros_rebaixamento.tex (relatorio_rebaixamento.pdf)
\newcommand{\RbCovPfModos}{6{,}2} % numeros_rebaixamento.tex (relatorio_rebaixamento.pdf)
\newcommand{\RbDifModos}{1{,}58} % numeros_rebaixamento.tex (relatorio_rebaixamento.pdf)
\newcommand{\RbGanhoModos}{1{,}2} % numeros_rebaixamento.tex (relatorio_rebaixamento.pdf)
\newcommand{\RbKNPar}{2000} % numeros_rebaixamento.tex (relatorio_rebaixamento.pdf)
\newcommand{\RbKPfFixo}{13{,}80} % numeros_rebaixamento.tex (relatorio_rebaixamento.pdf)
\newcommand{\RbKPfAleat}{15{,}70} % numeros_rebaixamento.tex (relatorio_rebaixamento.pdf)
\newcommand{\RbKMcNemar}{= 0{,}002} % numeros_rebaixamento.tex (relatorio_rebaixamento.pdf)
\newcommand{\SoBuracoAbaixo}{0} % numeros_solos.tex (relatorio_solos.pdf)
\newcommand{\SoBuracoAcima}{15} % numeros_solos.tex (relatorio_solos.pdf)
\newcommand{\SoBuracoFaixaAbaixo}{$[0{,}9975;\, 1)$} % numeros_solos.tex (relatorio_solos.pdf)
\newcommand{\SoBuracoFaixaAcima}{$[1;\, 1{,}0025)$} % numeros_solos.tex (relatorio_solos.pdf)
\newcommand{\SoBuracoPontos}{1{,}0} % numeros_solos.tex (relatorio_solos.pdf)
\newcommand{\SoBuracoFalta}{10} % numeros_solos.tex (relatorio_solos.pdf)
\newcommand{\SoNCinquenta}{1000} % numeros_solos.tex (relatorio_solos.pdf)
\newcommand{\ChSolo}{f50} % chuva_ai/ref48_piloto/mc_parametros.txt
\newcommand{\ChI}{31{,}05} % chuva_ai/ref48_piloto/mc_parametros.txt
\newcommand{\ChChuva}{24} % chuva_ai/ref48_piloto/mc_parametros.txt
\newcommand{\ChHw}{5} % chuva_ai/ref48_piloto/mc_parametros.txt
\newcommand{\ChTempoPasso}{6} % chuva_ai/ref48_piloto/mc_parametros.txt
\newcommand{\ChTempoFim}{48} % chuva_ai/ref48_piloto/mc_parametros.txt
\newcommand{\ChNTempos}{8} % chuva_ai/ref48_piloto/mc_parametros.txt
\newcommand{\ChCovC}{0{,}3} % chuva_ai/ref48_piloto/mc_parametros.txt
\newcommand{\ChCovPhi}{0{,}1} % chuva_ai/ref48_piloto/mc_parametros.txt
\newcommand{\ChCovK}{0{,}6} % chuva_ai/ref48_piloto/mc_parametros.txt
\newcommand{\ChLx}{20} % chuva_ai/ref48_piloto/mc_parametros.txt
\newcommand{\ChLy}{2} % chuva_ai/ref48_piloto/mc_parametros.txt
\newcommand{\ChKLM}{2000} % chuva_ai/ref48_piloto/mc_parametros.txt
\newcommand{\ChKLvar}{96{,}4} % chuva_ai/ref48_piloto/mc_parametros.txt
\newcommand{\ChMalhaT}{tri2} % chuva_ai/ref48_piloto/mc_parametros.txt
\newcommand{\ChMalhaF}{tri1} % chuva_ai/ref48_piloto/mc_parametros.txt
\newcommand{\ChNref}{4} % chuva_ai/ref48_piloto/mc_parametros.txt
\newcommand{\ChUt}{2} % chuva_ai/ref48_piloto/mc_parametros.txt
\newcommand{\ChSementeP}{2025} % chuva_ai/ref48_piloto/mc_parametros.txt
\newcommand{\ChSementeA}{1} % chuva_ai/ref48_ai/mc_parametros.txt
\newcommand{\ChBandaT}{24} % chuva_ai/ref48_ai/mc_parametros.txt
\newcommand{\ChEliteFrac}{10} % chuva_ai/ref48_ai/mc_parametros.txt
\newcommand{\ChNAlvoP}{200} % cluster/casos_chuva_ai.txt
\newcommand{\ChNAlvoA}{200} % cluster/casos_chuva_ai.txt
\newcommand{\ChDetMin}{1{,}1251} % chuva_ai/ref48_piloto/mc_0_0.csv (realizacao -1)
\newcommand{\ChDetT}{24} % chuva_ai/ref48_piloto/mc_0_0.csv (realizacao -1)
\newcommand{\ChForaP}{0} % chuva_ai/ref48_piloto/mc_*.csv (sem fator ou sem algum instante)
\newcommand{\ChForaA}{0} % chuva_ai/ref48_ai/mc_*.csv (sem fator ou sem algum instante)
\newcommand{\ChPIC}{exato} % IC do piloto (13 falhas)
\newcommand{\ChNP}{200} % chuva_ai/ref48_piloto/mc_*.csv
\newcommand{\ChPNf}{13} % chuva_ai/ref48_piloto/mc_*.csv (FS < 1 em algum instante)
\newcommand{\ChPPf}{6{,}50} % chuva_ai/ref48_piloto/mc_*.csv
\newcommand{\ChPLo}{3{,}5} % chuva_ai/ref48_piloto/mc_*.csv (IC exato)
\newcommand{\ChPHi}{10{,}9} % chuva_ai/ref48_piloto/mc_*.csv (IC exato)
\newcommand{\ChPCov}{26{,}9} % chuva_ai/ref48_piloto/mc_*.csv (s_q / (raiz(N) P_f))
\newcommand{\ChMinNoDet}{199} % chuva_ai/ref48_piloto/mc_*.csv (realizacoes com o minimo no instante do deterministico)
\newcommand{\ChFalhaFora}{0} % chuva_ai/ref48_piloto/mc_*.csv e chuva_ai/ref48_ai/mc_*.csv (FS < 1 noutro instante e >= 1 em 24 h)
\newcommand{\ChFalhaForaFrase}{nenhuma realização das duas amostras falha} % chuva_ai/ref48_piloto/mc_*.csv e chuva_ai/ref48_ai/mc_*.csv
\newcommand{\ChNA}{200} % chuva_ai/ref48_ai/mc_*.csv
\newcommand{\ChANf}{77} % chuva_ai/ref48_ai/mc_*.csv (falhas sorteadas)
\newcommand{\ChAFrac}{38} % chuva_ai/ref48_ai/mc_*.csv
\newcommand{\ChAPf}{6{,}22} % chuva_ai/ref48_ai/mc_*.csv (soma(w I) / N)
\newcommand{\ChAPfArred}{6} % chuva_ai/ref48_ai/mc_*.csv (arredondada)
\newcommand{\ChACov}{13{,}6} % chuva_ai/ref48_ai/mc_*.csv (s_q / (raiz(N) P_f))
\newcommand{\ChALo}{4{,}6} % chuva_ai/ref48_ai/mc_*.csv (P_f (1 - 1,96 CoV))
\newcommand{\ChAHi}{7{,}9} % chuva_ai/ref48_ai/mc_*.csv (P_f (1 + 1,96 CoV))
\newcommand{\ChAPesoMedio}{0{,}939} % chuva_ai/ref48_ai/mc_*.csv
\newcommand{\ChAPesoSe}{0{,}117} % chuva_ai/ref48_ai/mc_*.csv (desvio da media)
\newcommand{\ChAPesoCompat}{compatível com 1} % chuva_ai/ref48_ai/mc_*.csv (|media - 1| / desvio = 0.52)
\newcommand{\ChAEss}{49} % chuva_ai/ref48_ai/mc_*.csv ((soma w)^2 / soma w^2)
\newcommand{\ChAEssF}{43} % chuva_ai/ref48_ai/mc_*.csv (ESS das falhas: (soma q)^2 / soma q^2)
\newcommand{\ChPEssF}{13} % chuva_ai/ref48_piloto/mc_*.csv (ESS das falhas: no direto, n_f)
\newcommand{\ChPParcela}{8} % chuva_ai/ref48_piloto/mc_*.csv (maior parcela: no direto, 1 / n_f)
\newcommand{\AlertaParcela}{10} % ../README.md (criterios gerais)
\newcommand{\ChAPesoFalha}{0{,}16} % chuva_ai/ref48_ai/mc_*.csv (peso medio nas falhas)
\newcommand{\ChAPesoFalhaMax}{0{,}67} % chuva_ai/ref48_ai/mc_*.csv (maior peso nas falhas)
\newcommand{\ChAPesoMax}{9{,}8} % chuva_ai/ref48_ai/mc_*.csv (maior peso)
\newcommand{\ChAParcela}{5} % chuva_ai/ref48_ai/mc_*.csv (maior w I / soma w I)
\newcommand{\ChDif}{\ensuremath{-}0{,}14} % chuva_ai/ref48_ai/mc_*.csv e chuva_ai/ref48_piloto/mc_*.csv (desvios)
\newcommand{\ChDifPontos}{\ensuremath{-}0{,}28} % chuva_ai/ref48_ai/mc_*.csv e chuva_ai/ref48_piloto/mc_*.csv
\newcommand{\ChGanho}{4{,}1} % chuva_ai/ref48_ai/mc_*.csv ((CoV do direto com N e P_f da AI / CoV da AI)^2)
\newcommand{\ChRazao}{0{,}24} % chuva_ai/ref48_ai/mc_*.csv (1 / ganho)
\newcommand{\ChCovMesmoN}{27{,}5} % raiz((1 - P_f) / (N P_f)) com N e P_f da AI
\newcommand{\ChNParaCov}{820} % (1 - P_f) / (P_f CoV^2) com os da AI
\newcommand{\ChVies}{sem sinal de viés} % chuva_ai/ref48_ai/mc_*.csv
\newcommand{\ChTSeis}{6} % chuva_ai/ref48_piloto/mc_parametros.txt (TEMPOS)
\newcommand{\ChPfASeis}{0{,}00} % chuva_ai/ref48_ai/mc_*.csv (P_f em 6 h)
\newcommand{\ChCovASeis}{--} % chuva_ai/ref48_ai/mc_*.csv (CoV em 6 h)
\newcommand{\ChNfPSeis}{0} % chuva_ai/ref48_piloto/mc_*.csv (FS < 1 em 6 h)
\newcommand{\ChFalhasPSeis}{nenhuma falha} % chuva_ai/ref48_piloto/mc_*.csv (FS < 1 em 6 h)
\newcommand{\ChTDoze}{12} % chuva_ai/ref48_piloto/mc_parametros.txt (TEMPOS)
\newcommand{\ChPfADoze}{0{,}00} % chuva_ai/ref48_ai/mc_*.csv (P_f em 12 h)
\newcommand{\ChCovADoze}{--} % chuva_ai/ref48_ai/mc_*.csv (CoV em 12 h)
\newcommand{\ChNfPDoze}{0} % chuva_ai/ref48_piloto/mc_*.csv (FS < 1 em 12 h)
\newcommand{\ChFalhasPDoze}{nenhuma falha} % chuva_ai/ref48_piloto/mc_*.csv (FS < 1 em 12 h)
\newcommand{\ChTDezoito}{18} % chuva_ai/ref48_piloto/mc_parametros.txt (TEMPOS)
\newcommand{\ChPfADezoito}{0{,}014} % chuva_ai/ref48_ai/mc_*.csv (P_f em 18 h)
\newcommand{\ChCovADezoito}{61} % chuva_ai/ref48_ai/mc_*.csv (CoV em 18 h)
\newcommand{\ChNfPDezoito}{0} % chuva_ai/ref48_piloto/mc_*.csv (FS < 1 em 18 h)
\newcommand{\ChFalhasPDezoito}{nenhuma falha} % chuva_ai/ref48_piloto/mc_*.csv (FS < 1 em 18 h)
\newcommand{\ChTVinteQuatro}{24} % chuva_ai/ref48_piloto/mc_parametros.txt (TEMPOS)
\newcommand{\ChPfAVinteQuatro}{6{,}22} % chuva_ai/ref48_ai/mc_*.csv (P_f em 24 h)
\newcommand{\ChCovAVinteQuatro}{14} % chuva_ai/ref48_ai/mc_*.csv (CoV em 24 h)
\newcommand{\ChNfPVinteQuatro}{13} % chuva_ai/ref48_piloto/mc_*.csv (FS < 1 em 24 h)
\newcommand{\ChFalhasPVinteQuatro}{13 falhas} % chuva_ai/ref48_piloto/mc_*.csv (FS < 1 em 24 h)
\newcommand{\ChTTrinta}{30} % chuva_ai/ref48_piloto/mc_parametros.txt (TEMPOS)
\newcommand{\ChPfATrinta}{0{,}37} % chuva_ai/ref48_ai/mc_*.csv (P_f em 30 h)
\newcommand{\ChCovATrinta}{31} % chuva_ai/ref48_ai/mc_*.csv (CoV em 30 h)
\newcommand{\ChNfPTrinta}{1} % chuva_ai/ref48_piloto/mc_*.csv (FS < 1 em 30 h)
\newcommand{\ChFalhasPTrinta}{1 falha} % chuva_ai/ref48_piloto/mc_*.csv (FS < 1 em 30 h)
\newcommand{\ChTTrintaSeis}{36} % chuva_ai/ref48_piloto/mc_parametros.txt (TEMPOS)
\newcommand{\ChPfATrintaSeis}{0{,}035} % chuva_ai/ref48_ai/mc_*.csv (P_f em 36 h)
\newcommand{\ChCovATrintaSeis}{64} % chuva_ai/ref48_ai/mc_*.csv (CoV em 36 h)
\newcommand{\ChNfPTrintaSeis}{0} % chuva_ai/ref48_piloto/mc_*.csv (FS < 1 em 36 h)
\newcommand{\ChFalhasPTrintaSeis}{nenhuma falha} % chuva_ai/ref48_piloto/mc_*.csv (FS < 1 em 36 h)
\newcommand{\ChTQuarentaDois}{42} % chuva_ai/ref48_piloto/mc_parametros.txt (TEMPOS)
\newcommand{\ChPfAQuarentaDois}{0{,}011} % chuva_ai/ref48_ai/mc_*.csv (P_f em 42 h)
\newcommand{\ChCovAQuarentaDois}{73} % chuva_ai/ref48_ai/mc_*.csv (CoV em 42 h)
\newcommand{\ChNfPQuarentaDois}{0} % chuva_ai/ref48_piloto/mc_*.csv (FS < 1 em 42 h)
\newcommand{\ChFalhasPQuarentaDois}{nenhuma falha} % chuva_ai/ref48_piloto/mc_*.csv (FS < 1 em 42 h)
\newcommand{\ChTQuarentaOito}{48} % chuva_ai/ref48_piloto/mc_parametros.txt (TEMPOS)
\newcommand{\ChPfAQuarentaOito}{0{,}004} % chuva_ai/ref48_ai/mc_*.csv (P_f em 48 h)
\newcommand{\ChCovAQuarentaOito}{100} % chuva_ai/ref48_ai/mc_*.csv (CoV em 48 h)
\newcommand{\ChNfPQuarentaOito}{0} % chuva_ai/ref48_piloto/mc_*.csv (FS < 1 em 48 h)
\newcommand{\ChFalhasPQuarentaOito}{nenhuma falha} % chuva_ai/ref48_piloto/mc_*.csv (FS < 1 em 48 h)
\newcommand{\ChLinhasInstantes}{6 & 1{,}7777 & 1{,}7670 & 6{,}1 & 0 & 0 & 0 & -- \\
  12 & 1{,}4983 & 1{,}4788 & 6{,}7 & 0 & 0 & 0 & -- \\
  18 & 1{,}2751 & 1{,}2664 & 7{,}1 & 0 & 0 & 0{,}014 & 61 \\
  24 & 1{,}1251 & 1{,}1297 & 7{,}7 & 13 & 6{,}50 & 6{,}22 & 14 \\
  30 & 1{,}1973 & 1{,}1979 & 7{,}0 & 1 & 0{,}50 & 0{,}37 & 31 \\
  36 & 1{,}2501 & 1{,}2447 & 6{,}9 & 0 & 0 & 0{,}035 & 64 \\
  42 & 1{,}2894 & 1{,}2824 & 6{,}8 & 0 & 0 & 0{,}011 & 73 \\
  48 & 1{,}3225 & 1{,}3162 & 6{,}7 & 0 & 0 & 0{,}004 & 100 \\} % chuva_ai/ref48_piloto/mc_*.csv, chuva_ai/ref48_ai/mc_*.csv (por instante)
\newcommand{\ChZerosA}{6 e 12} % chuva_ai/ref48_ai/mc_*.csv (instantes sem falha sorteada)
\newcommand{\ChZeroHi}{1{,}8} % IC exato com 0 falhas em 200
\newcommand{\ChDeslocBlocos}{200} % chuva_ai/ref48_ai/mc_desloc_*.txt
\newcommand{\ChDeslocIguais}{o mesmo} % chuva_ai/ref48_ai/mc_desloc_*.txt
\newcommand{\ChElite}{20} % chuva_ai/ref48_ai/mc_desloc_*.txt
\newcommand{\ChNivel}{1{,}0079} % chuva_ai/ref48_ai/mc_desloc_*.txt (FS do conjunto elite)
\newcommand{\ChPilotoFalhas}{13} % chuva_ai/ref48_ai/mc_desloc_*.txt
\newcommand{\ChSC}{\ensuremath{-}1{,}26} % chuva_ai/ref48_ai/mc_desloc_*.txt
\newcommand{\ChSPhi}{\ensuremath{-}0{,}80} % chuva_ai/ref48_ai/mc_desloc_*.txt
\newcommand{\ChSK}{\ensuremath{-}0{,}13} % chuva_ai/ref48_ai/mc_desloc_*.txt
\newcommand{\ChMuDois}{2{,}25} % chuva_ai/ref48_ai/mc_desloc_*.txt
\newcommand{\ChVarPeso}{8{,}5} % exp(|mu|^2) - 1
\newcommand{\ChProjMedia}{\ensuremath{\ensuremath{-}1{,}23;\ \ensuremath{-}0{,}76;\ \ensuremath{-}0{,}28}} % chuva_ai/ref48_ai/mc_*.csv (media de proj_* sorteadas)
\newcommand{\ChProjSe}{0{,}07} % 1 / raiz(N)
\newcommand{\ChProjMaxDif}{2{,}12} % chuva_ai/ref48_ai/mc_*.csv e chuva_ai/ref48_ai/mc_desloc_*.txt (desvios, o maior dos tres campos)
\newcommand{\ChProjMaxCampo}{$k$} % chuva_ai/ref48_ai/mc_*.csv (o campo da maior diferenca)
\newcommand{\ChProjOutrasDif}{0{,}46 e 0{,}63} % chuva_ai/ref48_ai/mc_*.csv (os outros dois campos)
\newcommand{\ChBeta}{1{,}54} % formula, P_f da AI
\newcommand{\ChMdois}{3{,}88} % formula, P_f da AI (m = phi(beta) / P_f)
\newcommand{\ChGanhoLinear}{8{,}2} % formula, P_f da AI (exp(m^2) Phi(-beta - m))
\newcommand{\ChGanhoLinearMu}{7{,}8} % formula, P_f da AI e |mu| do deslocamento usado
\newcommand{\ChFracGanho}{50} % ganho medido / ganho linear
\newcommand{\ChProcessos}{48} % chuva_ai/logs/p*.txt
\newcommand{\ChNos}{4} % chuva_ai/logs/p*.txt (nos)
\newcommand{\ChInicio}{8h47} % chuva_ai/logs/p*.txt
\newcommand{\ChFim}{10h21} % chuva_ai/logs/p*.txt
\newcommand{\ChRelogio}{1~h 34~min} % chuva_ai/logs/p*.txt
\newcommand{\ChCustoP}{531} % chuva_ai/logs/p*.txt (s por realizacao, piloto)
\newcommand{\ChCustoA}{575} % chuva_ai/logs/p*.txt (s por realizacao, AI)
\newcommand{\ChCpu}{61} % chuva_ai/logs/p*.txt (soma dos tempos, h)
\newcommand{\ChVelhoN}{8} % Projects/SlopeRandomFieldRefinement/mc/gim_I/piloto_srm_I31.05_0_7.csv
\newcommand{\ChVelhoPrimeira}{0} % Projects/SlopeRandomFieldRefinement/mc/gim_I/piloto_srm_I31.05_0_7.csv
\newcommand{\ChVelhoUltima}{7} % Projects/SlopeRandomFieldRefinement/mc/gim_I/piloto_srm_I31.05_0_7.csv
\newcommand{\ChVelhoMedia}{1{,}1431} % Projects/SlopeRandomFieldRefinement/mc/gim_I/piloto_srm_I31.05_0_7.csv
\newcommand{\ChVelhoCov}{5{,}6} % Projects/SlopeRandomFieldRefinement/mc/gim_I/piloto_srm_I31.05_0_7.csv
\newcommand{\ChVelhoMin}{1{,}0339} % Projects/SlopeRandomFieldRefinement/mc/gim_I/piloto_srm_I31.05_0_7.csv
\newcommand{\ChVelhoLogn}{0{,}90} % Projects/SlopeRandomFieldRefinement/mc/gim_I/piloto_srm_I31.05_0_7.csv (ajuste log-normal)
\newcommand{\ChVelhoNf}{0} % Projects/SlopeRandomFieldRefinement/mc/gim_I/piloto_srm_I31.05_0_7.csv (FS < 1)
\newcommand{\ChVelhoHi}{37} % IC exato com 0 falhas em 8
\newcommand{\ChVelhoDifMax}{0{,}0015} % chuva_ai/ref48_piloto/mc_*.csv e Projects/SlopeRandomFieldRefinement/mc/gim_I/piloto_srm_I31.05_0_7.csv (maior |diferenca| de FS(24 h))
\newcommand{\ChOitoMedia}{1{,}1429} % chuva_ai/ref48_piloto/mc_*.csv (realizacoes 0..7, FS(24 h))
\newcommand{\ChOitoCov}{5{,}6} % chuva_ai/ref48_piloto/mc_*.csv (realizacoes 0..7)
\newcommand{\ChDuzMedia}{1{,}1297} % chuva_ai/ref48_piloto/mc_*.csv (FS(24 h))
\newcommand{\ChDuzCov}{7{,}7} % chuva_ai/ref48_piloto/mc_*.csv (FS(24 h))
\newcommand{\ChDuzLogn}{6{,}03} % chuva_ai/ref48_piloto/mc_*.csv (ajuste log-normal de FS(24 h))
\newcommand{\ChCruzA}{1{,}7} % log-normal: media de ln FS das 200, desvio das 8
\newcommand{\ChCruzB}{4{,}2} % log-normal: media de ln FS das 8, desvio das 200
\newcommand{\ChPfEsperada}{1} % cluster/README.md (pilotos do relatorio_chuva.tex)
\newcommand{\ChCovEspLin}{14} % cluster/README.md
\newcommand{\ChCovEspMax}{25} % cluster/README.md
\newcommand{\ChCovAlvo}{10} % COV_ALVO do script
\newcommand{\ChNAlvoCov}{1500} % (1 - P_f) / (P_f CoV^2) com a P_f da AI, em centenas
\newcommand{\ChCpuAlvo}{221} % N (em centenas) x custo do piloto no fermi
\newcommand{\ChRelogioAlvo}{4{,}6} % N (em centenas) x custo do piloto / processos
\newcommand{\ChoBloco}{50} % rodar.sh
\newcommand{\ChoC}{10} % campanha/r05_direto/mc_parametros.txt
\newcommand{\ChoPhi}{30} % campanha/r05_direto/mc_parametros.txt
\newcommand{\ChoGamma}{20} % campanha/r05_direto/mc_parametros.txt
\newcommand{\ChoE}{20} % campanha/r05_direto/mc_parametros.txt (MPa)
\newcommand{\ChoNu}{0{,}3} % campanha/r05_direto/mc_parametros.txt
\newcommand{\ChoCovC}{0{,}3} % campanha/r05_direto/mc_parametros.txt
\newcommand{\ChoCovPhi}{0{,}2} % campanha/r05_direto/mc_parametros.txt
\newcommand{\ChoLh}{20} % campanha/r05_direto/mc_parametros.txt
\newcommand{\ChoLv}{2} % campanha/r05_direto/mc_parametros.txt
\newcommand{\ChoKLM}{2000} % campanha/r05_direto/mc_parametros.txt
\newcommand{\ChoKLvar}{98{,}8} % campanha/r05_direto/mc_parametros.txt
\newcommand{\ChoSemente}{2010} % campanha/r05_direto/mc_parametros.txt
\newcommand{\ChoMalha}{tri2} % campanha/r05_direto/mc_parametros.txt
\newcommand{\ChoNref}{4} % campanha/r05_direto/mc_parametros.txt
\newcommand{\ChoBishop}{1{,}204} % Projects/Cho2010/README.md (Cho 2010)
\newcommand{\ChoPfCinco}{1{,}71} % Projects/Cho2010/README.md (Cho 2010): Tabela 5
\newcommand{\ChoPfSete}{0{,}4} % Projects/Cho2010/README.md (Cho 2010): Fig. 9 (aproximado)
\newcommand{\ChoNmc}{50\,000} % Projects/Cho2010/README.md (Cho 2010)
\newcommand{\ChoBetaCinco}{2{,}12} % -Phi^-1(P_f do Cho)
\newcommand{\ChoRCinco}{\ensuremath{-}0{,}5} % casos.txt
\newcommand{\ChoRSete}{\ensuremath{-}0{,}7} % casos.txt
\newcommand{\ChoRaiz}{0{,}866} % raiz(1 - r^2)
\newcommand{\ChoEliteFrac}{10} % cho2010.cpp (padrao de elite)
\newcommand{\ChoPilotoN}{1000} % casos.txt (piloto_n)
\newcommand{\ChoEliteN}{100} % elite x piloto_n
\newcommand{\ChoDetTriUm}{1{,}2097} % det/srm_fator_tri1_nref_4_0_-1.csv
\newcommand{\ChoEqTriUm}{4298} % det/srm_fator_tri1_nref_4_0_-1.csv (equacoes no ultimo nivel)
\newcommand{\ChoNivTriUm}{4} % det/srm_fator_tri1_nref_4_0_-1.csv (niveis)
\newcommand{\ChoDifTriUm}{+0{,}48} % det/srm_fator_tri1_nref_4_0_-1.csv e Projects/Cho2010/README.md (Cho 2010)
\newcommand{\ChoDetTriDoisTres}{1{,}2101} % det/srm_fator_tri2_nref_3_0_-1.csv
\newcommand{\ChoEqTriDoisTres}{4492} % det/srm_fator_tri2_nref_3_0_-1.csv (equacoes no ultimo nivel)
\newcommand{\ChoNivTriDoisTres}{3} % det/srm_fator_tri2_nref_3_0_-1.csv (niveis)
\newcommand{\ChoDifTriDoisTres}{+0{,}51} % det/srm_fator_tri2_nref_3_0_-1.csv e Projects/Cho2010/README.md (Cho 2010)
\newcommand{\ChoDetTriDois}{1{,}2047} % det/srm_fator_tri2_nref_4_0_-1.csv
\newcommand{\ChoEqTriDois}{7744} % det/srm_fator_tri2_nref_4_0_-1.csv (equacoes no ultimo nivel)
\newcommand{\ChoNivTriDois}{4} % det/srm_fator_tri2_nref_4_0_-1.csv (niveis)
\newcommand{\ChoDifTriDois}{+0{,}06} % det/srm_fator_tri2_nref_4_0_-1.csv e Projects/Cho2010/README.md (Cho 2010)
\newcommand{\ChoDetGim}{1{,}7899} % det/gim_0_-1.csv
\newcommand{\ChoEqGim}{8606} % det/gim_0_-1.csv (equacoes no ultimo nivel)
\newcommand{\ChoNivGim}{4} % det/gim_0_-1.csv (niveis)
\newcommand{\ChoUltNivTriDois}{\ensuremath{-}0{,}45} % det/srm_fator_tri2_nref_4 (niveis)
\newcommand{\ChoUltNivTriUm}{\ensuremath{-}0{,}81} % det/srm_fator_tri1_nref_4 (niveis)
\newcommand{\ChoProcessos}{14} % campanha/logs/pp*.txt
\newcommand{\ChoInicio}{9h12} % campanha/logs/pp*.txt
\newcommand{\ChoHoraDados}{11h32} % campanha/*/mc_*.csv (o mais novo)
\newcommand{\ChoCustoFalha}{393} % campanha/r0*_direto/mc_*.log (s, falhas)
\newcommand{\ChoCustoOutras}{94} % campanha/r0*_direto/mc_*.log (s, as outras)
\newcommand{\ChoCustoMedio}{97} % campanha/r0*_direto/mc_*.log (s)
\newcommand{\ChoCustoRazao}{4{,}2} % campanha/r0*_direto/mc_*.log
\newcommand{\ChoFeitas}{1213} % campanha/*/mc_*.csv
\newcommand{\ChoTotal}{7000} % casos.txt
\newcommand{\PcParou}{11h33} % ../README.md
\newcommand{\ChoFimDireto}{--} % fila simulada (estimativa)
\newcommand{\ChoInicioAI}{--} % fila simulada (estimativa)
\newcommand{\ChoFimMin}{--} % fila simulada (limite de baixo)
\newcommand{\ChoEmCurso}{14} % campanha/logs/pp*.txt (blocos em curso)
\newcommand{\ChoCincoEmCurso}{7} % campanha/logs/pp*.txt (blocos de r05_direto em curso, uma realizacao cada)
\newcommand{\ChoSeteEmCurso}{7} % campanha/logs/pp*.txt (blocos de r07_direto em curso, uma realizacao cada)
\newcommand{\ChoCincoIC}{exato} % campanha/r05_direto/mc_*.csv (10 falhas)
\newcommand{\ChoCincoN}{599} % campanha/r05_direto/mc_*.csv
\newcommand{\ChoCincoAlvo}{4000} % casos.txt
\newcommand{\ChoCincoNf}{10} % campanha/r05_direto/mc_*.csv (FS < 1)
\newcommand{\ChoCincoFalhas}{10 falhas} % campanha/r05_direto/mc_*.csv (FS < 1)
\newcommand{\ChoCincoPf}{1{,}67} % campanha/r05_direto/mc_*.csv
\newcommand{\ChoCincoLo}{0{,}80} % campanha/r05_direto/mc_*.csv (IC exato)
\newcommand{\ChoCincoHi}{3{,}05} % campanha/r05_direto/mc_*.csv (IC exato)
\newcommand{\ChoCincoNivDif}{0{,}0049} % campanha/r05_direto/mc_*.csv (FS do nivel 3 - FS do nivel 4, media)
\newcommand{\ChoCincoNivDp}{0{,}0021} % campanha/r05_direto/mc_*.csv (o mesmo, desvio)
\newcommand{\ChoCincoNivNf}{10} % campanha/r05_direto/mc_*.csv (FS < 1 no penultimo nivel)
\newcommand{\ChoCincoNivTrocas}{0} % campanha/r05_direto/mc_*.csv (realizacoes que trocam de lado entre os dois niveis)
\newcommand{\ChoCincoCov}{31} % campanha/r05_direto/mc_*.csv (s_q / (raiz(N) P_f))
\newcommand{\ChoCincoMedia}{1{,}1885} % campanha/r05_direto/mc_*.csv
\newcommand{\ChoCincoCovFS}{8{,}1} % campanha/r05_direto/mc_*.csv
\newcommand{\ChoCincoLogn}{1{,}80} % campanha/r05_direto/mc_*.csv (ajuste log-normal)
\newcommand{\ChoCincoDentro}{dentro} % campanha/r05_direto/mc_*.csv e Projects/Cho2010/README.md (Cho 2010)
\newcommand{\ChoCincoFracCurso}{7} % P_f t_falha / (P_f t_falha + (1 - P_f) t_outras)
\newcommand{\ChoCincoFalhasCurso}{0{,}46} % realizacoes em curso x fracao esperada de falhas
\newcommand{\ChoAICincoAlvo}{1000} % casos.txt
\newcommand{\ChoAICincoSemente}{1} % casos.txt
\newcommand{\ChoAICincoElite}{--} % a AI ainda nao gravou o deslocamento
\newcommand{\ChoAICincoNivel}{--} % a AI ainda nao gravou o deslocamento
\newcommand{\ChoAICincoPilFalhas}{--} % a AI ainda nao gravou o deslocamento
\newcommand{\ChoAICincoMuDois}{--} % a AI ainda nao gravou o deslocamento
\newcommand{\ChoAICincoS}{--} % a AI ainda nao gravou o deslocamento
\newcommand{\ChoAICincoIguais}{--} % a AI ainda nao gravou o deslocamento
\newcommand{\ChoAICincoN}{--} % a AI ainda nao tem realizacoes
\newcommand{\ChoAICincoNf}{--} % a AI ainda nao tem realizacoes
\newcommand{\ChoAICincoFalhas}{--} % a AI ainda nao tem realizacoes
\newcommand{\ChoAICincoPf}{--} % a AI ainda nao tem realizacoes
\newcommand{\ChoAICincoCov}{--} % a AI ainda nao tem realizacoes
\newcommand{\ChoAICincoLo}{--} % a AI ainda nao tem realizacoes
\newcommand{\ChoAICincoHi}{--} % a AI ainda nao tem realizacoes
\newcommand{\ChoAICincoDif}{--} % a AI ainda nao tem realizacoes
\newcommand{\ChoAICincoGanho}{--} % a AI ainda nao tem realizacoes
\newcommand{\ChoAICincoPeso}{--} % a AI ainda nao tem realizacoes
\newcommand{\ChoAICincoEss}{--} % a AI ainda nao tem realizacoes
\newcommand{\ChoGanhoLinCinco}{24} % formula, P_f da AI (ou do direto, sem a AI)
\newcommand{\ChoSeteIC}{exato} % campanha/r07_direto/mc_*.csv (2 falhas)
\newcommand{\ChoSeteN}{614} % campanha/r07_direto/mc_*.csv
\newcommand{\ChoSeteAlvo}{1000} % casos.txt
\newcommand{\ChoSeteNf}{2} % campanha/r07_direto/mc_*.csv (FS < 1)
\newcommand{\ChoSeteFalhas}{2 falhas} % campanha/r07_direto/mc_*.csv (FS < 1)
\newcommand{\ChoSetePf}{0{,}33} % campanha/r07_direto/mc_*.csv
\newcommand{\ChoSeteLo}{0{,}04} % campanha/r07_direto/mc_*.csv (IC exato)
\newcommand{\ChoSeteHi}{1{,}17} % campanha/r07_direto/mc_*.csv (IC exato)
\newcommand{\ChoSeteNivDif}{0{,}0049} % campanha/r07_direto/mc_*.csv (FS do nivel 3 - FS do nivel 4, media)
\newcommand{\ChoSeteNivDp}{0{,}0020} % campanha/r07_direto/mc_*.csv (o mesmo, desvio)
\newcommand{\ChoSeteNivNf}{1} % campanha/r07_direto/mc_*.csv (FS < 1 no penultimo nivel)
\newcommand{\ChoSeteNivTrocas}{1} % campanha/r07_direto/mc_*.csv (realizacoes que trocam de lado entre os dois niveis)
\newcommand{\ChoSeteCov}{71} % campanha/r07_direto/mc_*.csv (s_q / (raiz(N) P_f))
\newcommand{\ChoSeteMedia}{1{,}1933} % campanha/r07_direto/mc_*.csv
\newcommand{\ChoSeteCovFS}{6{,}7} % campanha/r07_direto/mc_*.csv
\newcommand{\ChoSeteLogn}{0{,}45} % campanha/r07_direto/mc_*.csv (ajuste log-normal)
\newcommand{\ChoSeteDentro}{dentro} % campanha/r07_direto/mc_*.csv e Projects/Cho2010/README.md (Cho 2010)
\newcommand{\ChoSeteFracCurso}{1} % P_f t_falha / (P_f t_falha + (1 - P_f) t_outras)
\newcommand{\ChoSeteFalhasCurso}{0{,}09} % realizacoes em curso x fracao esperada de falhas
\newcommand{\ChoAISeteAlvo}{1000} % casos.txt
\newcommand{\ChoAISeteSemente}{1} % casos.txt
\newcommand{\ChoAISeteElite}{--} % a AI ainda nao gravou o deslocamento
\newcommand{\ChoAISeteNivel}{--} % a AI ainda nao gravou o deslocamento
\newcommand{\ChoAISetePilFalhas}{--} % a AI ainda nao gravou o deslocamento
\newcommand{\ChoAISeteMuDois}{--} % a AI ainda nao gravou o deslocamento
\newcommand{\ChoAISeteS}{--} % a AI ainda nao gravou o deslocamento
\newcommand{\ChoAISeteIguais}{--} % a AI ainda nao gravou o deslocamento
\newcommand{\ChoAISeteN}{--} % a AI ainda nao tem realizacoes
\newcommand{\ChoAISeteNf}{--} % a AI ainda nao tem realizacoes
\newcommand{\ChoAISeteFalhas}{--} % a AI ainda nao tem realizacoes
\newcommand{\ChoAISetePf}{--} % a AI ainda nao tem realizacoes
\newcommand{\ChoAISeteCov}{--} % a AI ainda nao tem realizacoes
\newcommand{\ChoAISeteLo}{--} % a AI ainda nao tem realizacoes
\newcommand{\ChoAISeteHi}{--} % a AI ainda nao tem realizacoes
\newcommand{\ChoAISeteDif}{--} % a AI ainda nao tem realizacoes
\newcommand{\ChoAISeteGanho}{--} % a AI ainda nao tem realizacoes
\newcommand{\ChoAISetePeso}{--} % a AI ainda nao tem realizacoes
\newcommand{\ChoAISeteEss}{--} % a AI ainda nao tem realizacoes
\newcommand{\ChoGanhoLinSete}{100} % formula, P_f da AI (ou do direto, sem a AI)
\newcommand{\EtapaUmPar}{1000} % cluster/casos_etapa1.txt (r05_antigo)
\newcommand{\EtapaUmChuva}{600} % cluster/casos_etapa1.txt (chuva_antigo)
\newcommand{\EtapaUmFzero}{0{,}8} % cluster/casos_etapa1.txt
\newcommand{\ChoAntesN}{599} % campanha/r05_direto/mc_*.csv (realizacoes 0..999)
\newcommand{\ChoAntesNf}{10} % campanha/r05_direto/mc_*.csv (0..999, FS < 1)
\newcommand{\ChoAntesPerto}{8} % campanha/r05_direto/mc_*.csv (FS em [1; 1.02))
\newcommand{\ChoAntesFaixa}{$[1;\ 1{,}02)$} % FAIXAS do script
\newcommand{\FormTolGNovo}{0{,}001} % ../README.md (etapa 2)
\newcommand{\FormTolBNovo}{0{,}01} % ../README.md (etapa 2)
\newcommand{\FormDezM}{10} % form_M10_r05/form_parcial.json
\newcommand{\FormDezVar}{60{,}1} % form_M10_r05/it0_centro_parametros.txt
\newcommand{\FormDezAvalIt}{41} % 4M + 1
\newcommand{\FormDezIt}{5} % form_M10_r05/form_parcial.json (iteracoes completas)
\newcommand{\FormDezUltIt}{4} % form_M10_r05/form_parcial.json
\newcommand{\FormDezAval}{205} % form_M10_r05/form_parcial.json
\newcommand{\FormDezItAtual}{5} % form_M10_r05/it*_0_0.csv
\newcommand{\FormDezFeitas}{8} % form_M10_r05/it5_*_0_0.csv
\newcommand{\FormDezBeta}{2{,}13} % form_M10_r05/form_parcial.json
\newcommand{\FormDezG}{0{,}0004} % form_M10_r05/form_parcial.json (g da ultima iteracao)
\newcommand{\FormDezPf}{1{,}66} % form_M10_r05/form_parcial.json (Phi(-beta))
\newcommand{\FormDezBetaProx}{2{,}13} % form_M10_r05/form_parcial.json (proximo ponto do HL-RF)
\newcommand{\FormDezPfProx}{1{,}65} % form_M10_r05/form_parcial.json (Phi(-beta) do proximo ponto)
\newcommand{\FormDezAlfaC}{0{,}07} % form_M10_r05/form_parcial.json (soma dos alfa^2 de c')
\newcommand{\FormDezAlfaPhi}{0{,}93} % form_M10_r05/form_parcial.json (soma dos alfa^2 de phi')
\newcommand{\FormDezTop}{$\varphi'_{0}$ (0{,}74), $\varphi'_{2}$ (0{,}13), $c'_{0}$ (0{,}05)} % form_M10_r05/form_parcial.json (maiores alfa^2)
\newcommand{\FormDezConvergiu}{--} % form_M10_r05/form_parcial.json
\newcommand{\FormDezSobe}{sobe} % form_M10_r05/form_parcial.json
\newcommand{\FormDezDBeta}{0{,}005} % form_M10_r05/form_parcial.json (mudanca de beta na ultima iteracao)
\newcommand{\FormDezPasso}{0{,}11} % form_M10_r05/form_parcial.json (|u_k - u_k-1| na ultima iteracao)
\newcommand{\FormDezPassoMin}{0{,}10} % form_M10_r05/form_parcial.json (menor |du| depois da iteracao 2)
\newcommand{\FormDezCritNovo}{já passa} % form_M10_r05/form_parcial.json e ../README.md (|g| = 0.0004, mudanca de beta = 0.0049)
\newcommand{\FormDezParalelos}{4} % form_M10_r05/form_parcial.json
\newcommand{\FormDezItMax}{6} % form_M10_r05/form_parcial.json
\newcommand{\FormVinteM}{20} % form_M20_r05/form_parcial.json
\newcommand{\FormVinteVar}{72{,}1} % form_M20_r05/it0_centro_parametros.txt
\newcommand{\FormVinteAvalIt}{81} % 4M + 1
\newcommand{\FormVinteIt}{3} % form_M20_r05/form_parcial.json (iteracoes completas)
\newcommand{\FormVinteUltIt}{2} % form_M20_r05/form_parcial.json
\newcommand{\FormVinteAval}{243} % form_M20_r05/form_parcial.json
\newcommand{\FormVinteItAtual}{3} % form_M20_r05/it*_0_0.csv
\newcommand{\FormVinteFeitas}{41} % form_M20_r05/it3_*_0_0.csv
\newcommand{\FormVinteBeta}{2{,}01} % form_M20_r05/form_parcial.json
\newcommand{\FormVinteG}{0{,}0016} % form_M20_r05/form_parcial.json (g da ultima iteracao)
\newcommand{\FormVintePf}{2{,}24} % form_M20_r05/form_parcial.json (Phi(-beta))
\newcommand{\FormVinteBetaProx}{2{,}02} % form_M20_r05/form_parcial.json (proximo ponto do HL-RF)
\newcommand{\FormVintePfProx}{2{,}15} % form_M20_r05/form_parcial.json (Phi(-beta) do proximo ponto)
\newcommand{\FormVinteAlfaC}{0{,}06} % form_M20_r05/form_parcial.json (soma dos alfa^2 de c')
\newcommand{\FormVinteAlfaPhi}{0{,}94} % form_M20_r05/form_parcial.json (soma dos alfa^2 de phi')
\newcommand{\FormVinteTop}{$\varphi'_{0}$ (0{,}67), $\varphi'_{2}$ (0{,}13), $\varphi'_{12}$ (0{,}07)} % form_M20_r05/form_parcial.json (maiores alfa^2)
\newcommand{\FormVinteConvergiu}{--} % form_M20_r05/form_parcial.json
\newcommand{\FormVinteSobe}{sobe} % form_M20_r05/form_parcial.json
\newcommand{\FormVinteDBeta}{0{,}188} % form_M20_r05/form_parcial.json (mudanca de beta na ultima iteracao)
\newcommand{\FormVintePasso}{0{,}32} % form_M20_r05/form_parcial.json (|u_k - u_k-1| na ultima iteracao)
\newcommand{\FormVintePassoMin}{0{,}32} % form_M20_r05/form_parcial.json (menor |du| depois da iteracao 2)
\newcommand{\FormVinteCritNovo}{ainda não passa} % form_M20_r05/form_parcial.json e ../README.md (|g| = 0.0016, mudanca de beta = 0.1885)
\newcommand{\FormVinteParalelos}{4} % form_M20_r05/form_parcial.json
\newcommand{\FormVinteItMax}{5} % form_M20_r05/form_parcial.json
\newcommand{\FormH}{0{,}5} % form_M10_r05/form_parcial.json (h)
\newcommand{\FormTolG}{0{,}005} % form_kl.py
\newcommand{\FormTolP}{0{,}05} % form_kl.py
\newcommand{\FormCentroFS}{1{,}1686} % form_M10_r05/it0_centro_0_0.csv
\newcommand{\FormCentroC}{9{,}58} % form_M10_r05/it0_centro_0_0.csv (c_med)
\newcommand{\FormCentroPhi}{29{,}4} % form_M10_r05/it0_centro_0_0.csv (phi_med)
\newcommand{\FormGradC}{0{,}016} % form_M10_r05: gradiente em c'_0
\newcommand{\FormGradPhi}{0{,}062} % form_M10_r05: gradiente em phi'_0
\newcommand{\FormViaC}{0{,}052} % regra da cadeia: o efeito do xi de c'_0 pelo c'
\newcommand{\FormViaPhi}{\ensuremath{-}0{,}036} % regra da cadeia: o mesmo xi pelo phi' (r B)
\newcommand{\ChoBetaDireto}{2{,}13} % -Phi^-1(P_f do direto ate agora)
\newcommand{\FormCompara}{--} % os dois FORM ainda nao terminaram
\newcommand{\AIFormAlvo}{1000} % cadeia_form.sh
\newcommand{\AIFormSemente}{3} % cadeia_form.sh
\newcommand{\AIFormProc}{4} % cadeia_form.sh (processos no comeco)
\newcommand{\AIFormProcMais}{10} % cadeia_form.sh (processos depois do rodar.sh)
\newcommand{\AIFormMuDois}{4{,}54} % form_M10_r05/form_parcial.json (beta previsto ao quadrado)
\newcommand{\AIFormVarPeso}{93} % exp(beta^2) - 1, previsto
\newcommand{\AIFormN}{--} % a AI com o ponto de projeto ainda nao tem realizacoes
\newcommand{\AIFormNf}{--} % a AI com o ponto de projeto ainda nao tem realizacoes
\newcommand{\AIFormFalhas}{--} % a AI com o ponto de projeto ainda nao tem realizacoes
\newcommand{\AIFormPf}{--} % a AI com o ponto de projeto ainda nao tem realizacoes
\newcommand{\AIFormCov}{--} % a AI com o ponto de projeto ainda nao tem realizacoes
\newcommand{\AIFormLo}{--} % a AI com o ponto de projeto ainda nao tem realizacoes
\newcommand{\AIFormHi}{--} % a AI com o ponto de projeto ainda nao tem realizacoes
\newcommand{\AIFormDif}{--} % a AI com o ponto de projeto ainda nao tem realizacoes
\newcommand{\AIFormGanho}{--} % a AI com o ponto de projeto ainda nao tem realizacoes
\newcommand{\AIFormPeso}{--} % a AI com o ponto de projeto ainda nao tem realizacoes
\newcommand{\AIFormEss}{--} % a AI com o ponto de projeto ainda nao tem realizacoes
\newcommand{\RbVmdBanda}{13} % numeros_rebaixamento.tex (P_f e CoV citados)
\newcommand{\RbVmdModos}{17} % numeros_rebaixamento.tex (P_f e CoV citados)
\newcommand{\ChVmd}{59} % chuva_ai (P_f e CoV do piloto e da AI)
\newcommand{\BalLinhas}{rebaixamento & banda & \RbPfDireto{} (\RbNDireto) & \RbPfBanda{} (\RbNBanda) & \RbDifBanda & \RbVmdBanda & \RbGanhoBanda & \RbGanhoLinear \\
  rebaixamento & \RbKModos{} modos & \RbPfDireto{} (\RbNDireto) & \RbPfModos{} (\RbNModos) & \RbDifModos & \RbVmdModos & \RbGanhoModos & \RbGanhoLinear \\
  chuva (\ChSolo) & banda (elite) & \ChPPf{} (\ChNP) & \ChAPf{} (\ChNA) & \ChDif & \ChVmd & \ChGanho & \ChGanhoLinear \\
  Cho, $r = \ChoRCinco$ & banda (elite) & \ChoCincoPf{} (\ChoCincoN)$^\dagger$ & \multicolumn{4}{l}{não rodou no PC} & \ChoGanhoLinCinco \\
  Cho, $r = \ChoRSete$ & banda (elite) & \ChoSetePf{} (\ChoSeteN)$^\dagger$ & \multicolumn{4}{l}{não rodou no PC} & \ChoGanhoLinSete \\
  Cho, $r = \ChoRCinco$ & ponto de projeto & \ChoCincoPf{} (\ChoCincoN)$^\dagger$ & \multicolumn{4}{l}{não rodou no PC} & \ChoGanhoLinCinco \\} % as macros das linhas
\newcommand{\BalNCasos}{3} % casos completos do balanco
\newcommand{\BalMaxDif}{1{,}58} % o maior |diferenca| dos casos completos (desvios)
\newcommand{\BalVies}{todas abaixo de dois desvios} % casos completos do balanco
\newcommand{\BalGanhoMin}{2{,}7} % ganhos medidos na banda
\newcommand{\BalGanhoMax}{4{,}1} % ganhos medidos na banda
\newcommand{\BalFracMin}{50} % ganho medido / ganho linear, na banda
\newcommand{\BalFracMax}{69} % ganho medido / ganho linear, na banda
\newcommand{\BuracoFaixa}{$[0{,}99;\ 1)$} % FAIXAS do script
\newcommand{\BuracoCab}{$[0{,}98;\ 0{,}99)$ & $[0{,}99;\ 1)$ & $[1;\ 1{,}01)$ & $[1{,}01;\ 1{,}02)$} % FAIXAS do script
\newcommand{\ChoBuracoAbaixo}{0} % campanha/r05_direto/mc_*.csv (FS na faixa (0.99, 1.0))
\newcommand{\ChoBuracoViz}{4 e 5} % campanha/r05_direto/mc_*.csv (faixas (0.98, 0.99) e (1.0, 1.01))
\newcommand{\BuracoLinhas}{chuva: piloto (direto) & 200 & 1 & 2 & 7 & 2 \\
  chuva: AI (sem peso) & 200 & 7 & 6 & 7 & 11 \\
  Cho, $r = \ChoRCinco$: direto (parcial) & 599 & 4 & 0 & 5 & 3 \\
  Cho, $r = \ChoRSete$: direto (parcial) & 614 & 0 & 1 & 3 & 2 \\} % contagens de FS nas faixas [(0.98, 0.99), (0.99, 1.0), (1.0, 1.01), (1.01, 1.02)]
\newcommand{\FormLinhas}{0 & 0{,}1686 & 0{,}0000 & 0{,}1686 & 0{,}0000 \\
  1 & 0{,}0153 & 1{,}9196 & 0{,}0156 & 1{,}8183 \\
  2 & 0{,}0008 & 2{,}1156 & 0{,}0016 & 2{,}0068 \\
  3 & 0{,}0005 & 2{,}1247 &  &  \\
  4 & 0{,}0004 & 2{,}1296 &  &  \\} % form_M*_r05/form*.json (historico)
\newcommand{\ChoLinhas}{$r = \ChoRCinco$ & direto & \ChoCincoN{} de \ChoCincoAlvo & \ChoCincoNf & \ChoCincoPf & \ChoCincoLo--\ChoCincoHi & \ChoCincoCov & \ChoPfCinco \\
   & AI na banda & \multicolumn{6}{l}{não rodou no PC (vai para o fermi)} \\
   & AI no ponto de projeto & \multicolumn{6}{l}{não rodou no PC (vai para o fermi)} \\
  \addlinespace
  $r = \ChoRSete$ & direto & \ChoSeteN{} de \ChoSeteAlvo & \ChoSeteNf & \ChoSetePf & \ChoSeteLo--\ChoSeteHi & \ChoSeteCov & $\approx$\ChoPfSete \\
   & AI na banda & \multicolumn{6}{l}{não rodou no PC (vai para o fermi)} \\} % as macros das linhas
